% Dermoscopy results. Every number is produced by scripts in the repository; see results/SUMMARY.md.
\section{Results: dermoscopy}\label{sec:derm}

\subsection{Masking reverses with location (controlled)}
With a synthetic ruler whose overlap with the lesion is swept from 0 to 100\% (2,437 ISIC 2018 images,
DINOv2 ViT-B/14 at 518\,px, three seeds), ROI masking improves reversed-environment AUROC over ERM by
$+0.184$ [$+0.146, +0.224$] when the ruler lies outside the lesion and \emph{lowers} it by $-0.097$
[$-0.138, -0.056$], $-0.112$ [$-0.155, -0.068$] and $-0.129$ [$-0.170, -0.087$] at 50, 75 and 100\% overlap
(location interaction $+0.313$ [$+0.272, +0.353$]; Table~\ref{tab:synthetic_isic_dino518}). The sign flip is
reproduced at 224\,px ($+0.361$ at 0\%; $-0.095$ at 100\%) and under randomised ruler colour, opacity, ticks and
scale. At 0\% overlap the masked model's correlated and reversed AUROC coincide (0.849/0.849): the shortcut is
removed. At 100\% the masked model's correlated AUROC exceeds ERM's (0.961 vs 0.937): it uses the ruler \emph{more}.

\begin{table}
\caption{Synthetic ruler, ISIC 2018, DINOv2 ViT-B/14 @518: AUROC and reversed-test Δ vs ERM [95% CI]}
\label{tab:synthetic_isic_dino518}
\begin{tabular}{lllllllllll}
\toprule
Method & Needs & clean@0 & rev@0 & clean@50 & rev@50 & clean@100 & rev@100 & Δrev@0 & Δrev@50 & Δrev@100 \\
\midrule
ERM & — & 0.797 & 0.665 & 0.793 & 0.573 & 0.792 & 0.515 &  &  &  \\
ROI masking & ROI mask (train+test) & 0.849 & 0.849 & 0.829 & 0.476 & 0.825 & 0.386 & +0.184 [+0.146, +0.224] & -0.097 [-0.138, -0.056] & -0.129 [-0.170, -0.087] \\
Inpainting (oracle) & artifact pixels (train+test) & 0.810 & 0.806 & 0.809 & 0.783 & 0.807 & 0.759 & +0.141 [+0.120, +0.162] & +0.211 [+0.187, +0.234] & +0.243 [+0.218, +0.269] \\
Group-balanced & image-level A (train) & 0.814 & 0.802 & 0.815 & 0.792 & 0.806 & 0.796 & +0.137 [+0.109, +0.168] & +0.219 [+0.188, +0.252] & +0.281 [+0.248, +0.315] \\
DFR & image-level A (val) & 0.786 & 0.789 & 0.787 & 0.794 & 0.789 & 0.794 & +0.124 [+0.091, +0.160] & +0.221 [+0.184, +0.261] & +0.279 [+0.240, +0.320] \\
LEACE (paired, oracle renders) & paired renders (train) & 0.829 & 0.832 & 0.810 & 0.810 & 0.809 & 0.809 & +0.167 [+0.139, +0.197] & +0.237 [+0.212, +0.263] & +0.294 [+0.265, +0.325] \\
Repair consistency & artifact pixels (train) & 0.828 & 0.707 & 0.822 & 0.616 & 0.814 & 0.558 & +0.042 [+0.015, +0.069] & +0.043 [+0.014, +0.072] & +0.042 [+0.012, +0.073] \\
I2E (ours) & template library & 0.848 & 0.844 & 0.850 & 0.833 & 0.847 & 0.818 & +0.179 [+0.152, +0.206] & +0.260 [+0.229, +0.291] & +0.303 [+0.271, +0.336] \\
I2E + balanced (ours) & templates + A (train) & 0.841 & 0.842 & 0.843 & 0.840 & 0.844 & 0.842 & +0.177 [+0.151, +0.205] & +0.267 [+0.238, +0.298] & +0.327 [+0.293, +0.362] \\
I2E rank-1 (ablation) & template library & 0.811 & 0.799 & 0.809 & 0.770 & 0.806 & 0.730 & +0.134 [+0.115, +0.154] & +0.197 [+0.174, +0.220] & +0.215 [+0.190, +0.240] \\
Insertion augmentation & template library & 0.814 & 0.726 & 0.812 & 0.654 & 0.810 & 0.603 & +0.061 [+0.049, +0.074] & +0.081 [+0.068, +0.094] & +0.088 [+0.077, +0.099] \\
\bottomrule
\end{tabular}
\end{table}

Insert-then-erase in this table is the artifact-specific precursor of U-MtE: its insertion pairs use templates of the
artifact type (rulers) and it erases from unmasked features. U-MtE (Sec.~\ref{sec:i2e}) replaces the templates by
one generic overlay library for every modality and erases after masking; its dermoscopy results are in
Sec.~\ref{sec:solution}.

\subsection{Mechanism: retention, not occlusion or context}
Three explanations were tested separately. \emph{Occlusion}: with the ruler present but uncorrelated with the
label (50/50), masking helps at every overlap ($+0.036$, $+0.040$, $+0.039$ at 0/50/100\%); covering lesion
pixels costs nothing measurable, so the harm requires the correlation. \emph{Lost context}: dilating the mask at
50\% overlap \emph{increases} harm ($-0.097 \rightarrow -0.138 \rightarrow -0.154$ for margins 0/10/25\,px) as the
margin restores ruler pixels (retention $0.50 \rightarrow 0.73 \rightarrow 0.95$); at 100\% overlap, where retention
is already 1, harm is flat ($-0.129$ to $-0.144$) although the ruler's share of visible pixels falls threefold.
\emph{Amplified reliance}: masking doubles the same-head counterfactual sensitivity to the ruler
($|\Delta p|$ 0.235 $\rightarrow$ 0.491 at 100\%). Harm is largest for small lesions ($-0.364$ vs $+0.064$ for large
lesions at 100\%), where the ruler is a larger share of what survives the mask.

\subsection{Real hair (ISIC 2019), pre-registered}
Two traps share one hair-free group (at most 30 hair pixels): Trap~A contains hair lying mostly inside the lesion
($r\ge0.5$), Trap~B hair mostly outside ($r<0.1$); cells are matched to the hair-free source mix, and five seeds
of five-fold group-safe cross-validation give 172 reversed-test melanomas per trap. Masking helps in Trap~B
($+0.107$ [$+0.078, +0.135$]), harms in Trap~A ($-0.037$ [$-0.057, -0.019$], negative in all five seeds) and the
crossover is $+0.145$ [$+0.105, +0.181$]. The harm holds in BCN images ($-0.051$ [$-0.075, -0.027$]); in HAM it
is not distinguishable from zero ($-0.009$ [$-0.036, +0.017$]). Oracle hair inpainting helps in both traps
($+0.103$, $+0.096$). Group-balanced training and DFR, which need only an image-level hair label, give the largest
gains ($+0.225$/$+0.251$ in Trap~A, $+0.209$/$+0.215$ in Trap~B). Paired linear erasure recovers less than half of
that ($+0.086$/$+0.081$). The magnitude of the in-ROI harm depends on the definition of the comparison group: with
a looser hair-free group ($<0.1\%$ of pixels) and only heavily haired images, masking gives no benefit but no
measurable harm inside the lesion ($+0.012$ [$-0.042, +0.062$]) while still helping outside ($+0.178$); the
benefit shrinks to zero as more of the hair lies inside the lesion.

With the dermatology foundation model DermLIP, the controlled ruler sweep reverses as with DINOv2 ($+0.161$ at no
overlap, $-0.263$ to $-0.309$ at 25--100\%), but on real hair masking harms at both locations (Trap~A $-0.085$
[$-0.107, -0.063$], Trap~B $-0.088$): with this encoder masking removes disease context regardless of where the hair
lies, so no location gap appears.

\subsection{What the real-hair shortcut is made of}
Removing hair pixels at test time only closes 44--47\% of the ERM clean--reversed gap: most of what a frozen
encoder exploits in a real hair trap is not the hair pixels but what accompanies hair (hair-bearing images differ in
sex, age and body site). This bounds every pixel-level method, including oracle inpainting and U-MtE, and
explains why image-level-label methods, which remove everything correlated with the label, recover more.

\subsection{Reversed-environment gains can be gamed}
Erasing the artifact concept from unpaired image-level labels yields the largest reversed gain in the study
($+0.326$ in Trap~A) while \emph{inverting} the correlated/reversed ordering (0.618 vs 0.817). Post-hoc
prevalence-equalised recalibration \citep{kina2026prevalence} behaves the same way under the trap's 90/10 shift
($+0.427$; correlated 0.433 vs reversed 0.918). Both remove diagnosis-correlated signal or over-correct the
group prior rather than removing reliance on the artifact; we therefore read every reversed gain together with
clean AUROC and the correlated/reversed ordering.
